\section{Background and Threat Model}\label{sec:background}

\paragraph{Syscall monitoring on Kubernetes.}
A container is a group of processes that share a set of namespaces and a control group (cgroup); every action it takes
on the host, from opening a file to changing a namespace, is a system call. eBPF lets a verified program run inside the
kernel at a hook such as a tracepoint, keep state in maps shared with user space, and filter by cgroup before doing any
work. Two raw tracepoints, \texttt{sys\_enter} and \texttt{sys\_exit}, see every system call of a node; the latter also
sees its return value, which carries the error code of a failure and the number of bytes a read or write moved.
Production runtime tools build on the same hooks: Falco~\cite{falco} and Tracee~\cite{tracee} match events against
rules, and Tetragon~\cite{tetragon} attaches kprobes to kernel functions and enforces policies. Rules catch what their
authors anticipated; the question this paper studies is how far a model learnt only from a workload's own normal
behaviour can go beyond them while staying fast and quiet enough for a live cluster.

\paragraph{Threat model.}
The attacker already runs code inside an application container, or exploits the containerised service itself. On the
cluster they start a new process inside the container and perform syscall-level actions: namespace changes, ptrace,
unusual exec, changes of memory protection, use of the seccomp API. On the public datasets the exploit goes through
the application. We protect every replica of the 21 workloads of one application namespace; infrastructure pods, load
generators and control-plane nodes are not scored. We trust the kernel, the eBPF runtime, Tetragon and our worker, and
we assume the normal traffic used for learning and thresholds contains no attack.

Out of scope are (i) an adaptive attacker who uses only system-call sequences already present in the workload's
normal behaviour, (ii) attacks that only change the meaning of data inside a process without changing its system
behaviour, such as an SQL injection that is an ordinary login request at the syscall level, and (iii) attacks on the
host or on the cluster infrastructure. We measure how the system fails under (i) and (ii) rather than claim robustness
(\cref{sec:discussion}). Our cluster campaign covers MITRE ATT\&CK for Containers~\cite{attack_containers} only
partially: execution, privilege escalation or escape attempts and defense impairment through self-written fixtures;
the public datasets add initial access through 15 real vulnerabilities and nine container-breakout scenarios.
